% ATTEMPT_STATUS: unresolved
% RESEARCH_GENERATED: 2026-10-01; GPT-6 Astra Ultra; individual mathematical attempt
% TOP_PROBLEM: {"id": 145, "title": "Trace Reconstruction", "category_id": 2, "category": {"id": 2, "name": "combinatorics", "display_name": "Combinatorics", "description": "Counting problems, graph theory, discrete structures.", "slug": "combinatorics", "order_index": 2, "created_at": "2026-07-31T15:26:25.671Z"}, "status": "open"}
% FIRST_POSTED: 2026-10-01
% Imported from: unsolvedmath_additional_500.tex; UnsolvedMath ID 145
% !TeX program = lualatex
% Compile twice: lualatex 145.tex
% Self-contained: no external bibliography, images, or \input files.
% Numeric section labels and visible numbers are original UnsolvedMath IDs.
\documentclass[11pt,a4paper]{article}
\usepackage[margin=24mm,headheight=14pt]{geometry}
\usepackage{fontspec}
\setmainfont{Latin Modern Roman}
\setsansfont{Latin Modern Sans}
\setmonofont{Latin Modern Mono}
\usepackage{amsmath,amssymb,mathtools}
\usepackage{unicode-math}
\setmathfont{Latin Modern Math}
\usepackage{microtype}
\usepackage{enumitem,xurl,needspace}
\usepackage[dvipsnames]{xcolor}
\usepackage{hyperref}
\hypersetup{unicode=true,colorlinks=true,linkcolor=MidnightBlue,urlcolor=MidnightBlue,
 pdftitle={UnsolvedMath Problem 145},pdfauthor={Source-annotated compilation},bookmarksnumbered=true}
\usepackage{bookmark,tocloft,titlesec,fancyhdr}
\setlength{\cftsecnumwidth}{4.2em}
\cftsetpnumwidth{2.5em}
\cftsetrmarg{4em}
\titleformat{\section}{\Large\bfseries\raggedright}{\thesection}{0.7em}{}
\pagestyle{fancy}\fancyhf{}
\fancyhead[L]{\small\sffamily UnsolvedMath research attempt}
\fancyhead[R]{\small Original catalogue IDs}
\fancyfoot[C]{\thepage}
\setlength{\parindent}{0pt}
\setlength{\parskip}{5pt plus 1pt}
\setlength{\emergencystretch}{3em}
\setcounter{tocdepth}{1}\setcounter{secnumdepth}{1}
\setlist[itemize]{leftmargin=3em,itemsep=4pt,topsep=4pt,parsep=0pt}
\allowdisplaybreaks
\newcommand{\N}{\mathbb{N}}
\newcommand{\Z}{\mathbb{Z}}
\newcommand{\Q}{\mathbb{Q}}
\newcommand{\R}{\mathbb{R}}
\newcommand{\C}{\mathbb{C}}
\newcommand{\F}{\mathbb{F}}
\newcommand{\A}{\mathbb{A}}
\newcommand{\PP}{\mathbb{P}}
\newcommand{\E}{\mathbb{E}}
\newcommand{\eps}{\varepsilon}
\newcommand{\dd}{\,\mathrm d}
\newcommand{\sref}[2]{\hyperlink{source:#1:#2}{[S#2]}}
\newif\ifshowcatalogue
\showcataloguetrue
\newcommand{\cataloguescope}[1]{\ifshowcatalogue
 \paragraph{Statement recorded in the supplied source.}
 #1\fi}
\begin{document}
% UnsolvedMath ID 145; source catalogue code GREEN-057

\setcounter{section}{144}

\section{Worst-Case Binary Trace Reconstruction}\label{145}

{\small\sffamily UnsolvedMath ID 145 · Combinatorics}

\subsection{Definitions and mathematical statement}

\paragraph{Shared notation.}Unless a section says otherwise, $\N=\{1,2,\ldots\}$,
$\Z$ is the ring of integers, $\Q,\R,\C$ are the rational, real, and complex
fields, $[N]=\{1,\ldots,\lfloor N\rfloor\}$, and $\log$ is the natural
logarithm. For real functions of a parameter tending to infinity,
$f\ll g$ and $f=O(g)$ mean $|f|\leq Cg$ eventually for some fixed $C$;
$f\gg g$ reverses this comparison; $f=o(g)$ means $f/g\to0$;
$f\sim g$ means $f/g\to1$; and $f\asymp g$ means both order bounds.
A subscript on $O$ or $\ll$ specifies allowed constant dependence.
The expression $n^{o(1)}$ in an upper bound means growth smaller than
$n^\epsilon$ for each fixed $\epsilon>0$. Local definitions govern any
exception, finite-field convention, or use of the same symbol for another
object. An asymptotic claim is not automatically an effective algorithm.



A trace of $x\in\{0,1\}^n$ is the subsequence left after independently deleting each coordinate with probability $1/2$, retaining the order but not the deleted positions. Let $m(n)$ be the least sample size for which one reconstruction rule, knowing $n$, returns every possible $x$ correctly with probability at least $0.9$ from independent traces. This is a sample-complexity question, without an additional running-time bound. \sref{145}{1} \sref{145}{2}

\paragraph{Statement recorded in the supplied source.}
Given a string $x \in \{0, 1\}^n$, let $\tilde{x}$ be obtained by deleting bits independently at random with probability $\frac{1}{2}$. How many independent traces $\tilde{x}_1, \dots, \tilde{x}_m$ are needed to reconstruct $x$ with probability 0.9? \sref{145}{1}

\paragraph{Residual task reported by the archive.}

The embedded review (2026-08-17) records: Determine matching lower and upper bounds, including the optimal exponent at deletion probability 1/2. \sref{145}{1}

\subsection{Short English statement}

How many independent deletion-corrupted copies suffice to recover an arbitrary binary string reliably, even when computation time is unrestricted?

\subsection{Research attempt}

\paragraph{Provenance and scope.}
This individual attempt was generated by GPT-6 Astra Ultra on 1 October 2026. The method was to derive explicit partial results and test the first proposed extension adversarially. The original statement and sources below are retained. No claim of mathematical novelty or complete solution is made.

\paragraph{Formal target and context.}
The unknown binary string has known length $n$, deletion probability is one half, and the same algorithm must succeed with probability at least $0.9$ for every input. Computation time is unrestricted. The inspected Nazarov--Peres paper concerns a much stronger upper bound than the elementary one below and distinguishes the restricted model based only on coordinate means. We derive a simple information-theoretic lower bound from neighbouring isolated-one strings and an exact generating-function identity, keeping general reconstruction separate from mean-based reconstruction.

\paragraph{An elementary upper bound uses an intact trace.}
A trace equals the complete input with probability $2^{-n}$. Whenever a trace of length $n$ appears, it must be the original string. With $m$ independent traces, the probability of seeing none is $(1-2^{-n})^m\leq e^{-m2^{-n}}$. Thus $m\geq2^n\log10$ suffices. This is deliberately a weak baseline, not a claim about the best known rate. Its merit is that it is uniform in the string and gives a direct check on the success-probability convention.

\paragraph{Two nearby isolated-one inputs are hard to distinguish.}
Take $x=0^k1 0^\ell$ and $y=0^{k+1}1 0^{\ell-1}$, where $k+\ell+1=n$ and both $k,\ell$ are order $n$. Conditional on deletion of the one, both inputs produce the same distribution: a zero string of length $\operatorname{Bin}(n-1,1/2)$. Conditional on retention, the numbers of zeros before and after the one are independent, distributed as
\[
 (\operatorname{Bin}(k,1/2),\operatorname{Bin}(\ell,1/2))
 \quad\hbox{and}\quad
 (\operatorname{Bin}(k+1,1/2),\operatorname{Bin}(\ell-1,1/2)).
\]
Let $P_m$ denote the law of $\operatorname{Bin}(m,1/2)$. Since $P_{m+1}$ is the average of $P_m$ and its translate by one,
\[
 d_{\rm TV}(P_m,P_{m+1})
 \leq\tfrac12d_{\rm TV}(P_m,P_m+1)
 =\tfrac12\max_j P_m(j)\ll m^{-1/2}.
\]
For completeness, the equality follows by summing successive absolute differences of the unimodal binomial mass function: the total rise and fall is twice its maximum, and total variation divides by two. The final estimate is the central-binomial bound from Stirling's formula.

Total variation between product measures is at most the sum of the coordinate distances, by independent maximal couplings. Mixing with the common deleted-one component multiplies the remaining difference by one half. Therefore the single-trace laws $P_x,P_y$ satisfy $d_{\rm TV}(P_x,P_y)\leq C/\sqrt n$. For $m$ independent traces, the same coupling argument gives distance at most $Cm/\sqrt n$.

If a reconstruction rule succeeds with probability $0.9$ on both inputs, its output distinguishes them with average error at most $0.1$. Binary testing requires total variation at least $0.8$: for equal priors, optimal average success is $(1+d_{\rm TV})/2$. Hence $m\geq c\sqrt n$. This is a complete but nonoptimal lower bound; it does not use computational restrictions.

\paragraph{The mean polynomial and its limitation.}
Index the input bits by $0,\ldots,n-1$, and write a trace's bit at position $j$ as $\widetilde x_j$, padded by zeros beyond its length. Conditional on retaining bit $i$, its output position is binomial with parameters $i,1/2$. Consequently
\[
 \E\sum_{j\geq0}\widetilde x_jz^j
 =\frac12\sum_{i=0}^{n-1}x_i\left(\frac{1+z}{2}\right)^i.
\]
The affine change of variable is invertible, so exact knowledge of all means identifies the input polynomial uniquely. However finite samples only estimate means approximately. Inverting a polynomial transformation can amplify tiny errors, and injectivity alone gives no good sample bound. This is the exact point where a purely algebraic identification argument stops being a quantitative reconstruction theorem.

\paragraph{Verification and remaining gap.}
The two test strings have the same length and the one is internal, so no boundary convention is hidden in the binomial laws. The common all-zero mixture is treated explicitly rather than discarded. The lower bound concerns every reconstruction algorithm, whereas the generating identity describes only means. The established bounds are $c\sqrt n\leq m(n)\leq2^n\log10$ by elementary arguments, together with exact identifiability from means. They are not literature records. Determining the sharp sample growth still requires substantially stronger indistinguishability examples or reconstruction estimates beyond these calculations.

\subsection{Sources}

{\small

\begin{itemize}

\item[\hypertarget{source:145:1}{S1}] ULAM.AI, \emph{UnsolvedMath}, supplied \texttt{problems.json}, record 145 (GREEN-057), statement and background. The embedded literature-review date is 2026-08-17; that is the archive’s review date, not a fresh certification. Dataset landing page: \url{https://huggingface.co/datasets/ulamai/UnsolvedMath}.

\item[\hypertarget{source:145:2}{S2}] B. Green, 100 Open Problems, maintained author notes. \url{https://people.maths.ox.ac.uk/greenbj/papers/open-problems.pdf}. The author-hosted document was inspected on 27 September 2026; the archive’s local code is not a problem-number locator in this paper.

\item[E1] F. Nazarov and Y. Peres, \emph{Trace reconstruction with $\exp(O(n^{1/3}))$ samples}. \url{https://arxiv.org/abs/1612.03599}. Inspected 1 October 2026.

\end{itemize}

}

\label{end:145}
\end{document}
